\subsection{替换}

设 K 是一个交换域，E 是一个含单位元的结合 K-代数， \(\mathbf{x} = (x_i)_{i \in I}\) 是 E 中两两可交换的元素的一个族。设 \(\mathbf{B} = \mathbf{K}[(\mathbf{X}_i)_{i \in I}]\) ，并设 \(S_x\) 为所有非零 \(v \in B\) （使 \(v(\mathbf{x})\) 在 E 中可逆者）组成的集合。设 \(u \in B\) , \(v \in S_x\) 且 \(f = \frac{u}{v} \in \mathbf{K}((\mathbf{X}_i)_{i \in I})\) 。元素 \(u(\mathbf{x}) v(\mathbf{x})^{-1} = v(\mathbf{x})^{-1} u(\mathbf{x})\) 在 E 中有定义；此外，如果 \(u_1\) , \(v_1\) 是两个多项式，使得 \(f = \frac{u_1}{v_1}\) 且 \(v_1 \in S_x\) ，则 \(uv_1 = vu_1\) ，因此 \(u(\mathbf{x}) v_1(\mathbf{x}) = v(\mathbf{x}) u_1(\mathbf{x})\) ，从而

\[
u (\mathbf {x}) v (\mathbf {x}) ^ {- 1} = u _ {1} (\mathbf {x}) v _ {1} (\mathbf {x}) ^ {- 1}.
\]

设 \(f \in \mathbf{K}((\mathbf{X}_i)_{i \in I})\) 。如果至少存在一对 \((u, v)\) 使得 \(f = \frac{u}{v}\) 且 \(v \in S_x\) ，我们就说 \(x\) 在 \(f\) 中可代入；元素 \(u(x) v(x)^{-1}\) （它仅依赖于 \(f\) 与 \(x\) ）这时记作 \(f(x)\) 或 \(f((x_i))\) 或 \(f((x_i)_{i \in I})\) 。

命题 2. --- 设 K 是一个交换域，E 是一个含单位元的结合 K-代数，且 \(\mathbf{x} = (x_i)_{i \in I}\) 是 E 中两两可交换的元素的一个族。集合 \(\mathbf{S}_{\mathbf{x}}^{-1}\mathbf{B}\) 由所有 \(f \in \mathbf{K}((\mathbf{X}_i)_{i \in I})\) （ \(\mathbf{x}\) 在 \(f\) 中可代入者）组成，它是 \(\mathbf{K}((\mathbf{X}_i)_{i \in I})\) 的一个 K-子代数。映射 \(f \mapsto f(\mathbf{x})\) 是一个保单位同态 \(\varphi\) （ \(\mathbf{S}_{\mathbf{x}}^{-1}\mathbf{B}\) 到 E 中的）。像 \(\varphi(\mathbf{S}_{\mathbf{x}}^{-1}\mathbf{B})\) 是所有 \(yz^{-1}\) （其中 \(y\) 遍历含单位元子代数 \(\mathbf{K}[\mathbf{x}]_E\) （ E 的、由族 \(\mathbf{x}\) 生成的），而 \(z\) 遍历 \(\mathbf{K}[\mathbf{x}]_E\) 的所有可逆元素）组成的集合。

设 \(f_{1} = \frac{u_{1}}{v_{1}}\) , \(f_{2} = \frac{u_{2}}{v_{2}}\) 是 \(K((X_{i})_{i\in I})\) 的两个元素，满足 \(v_{1}, v_{2} \in S_{x}\) 。我们有 \(f_{1} + f_{2} = \frac{u_{1}v_{2} + u_{2}v_{1}}{v_{1}v_{2}}\) , \(f_{1}f_{2} = \frac{u_{1}u_{2}}{v_{1}v_{2}}\) 且 \(v_{1}, v_{2} \in S_{x}\) 。因此 \(S_{x}^{-1}B\) 是 \(K((X_{i})_{i\in I})\) 的一个 K-子代数。命题的其余部分是显然的。

推论。--- 设 L 是一个交换域，K 是 L 的一个子域， \(\mathbf{x} = (x_i)_{i \in I}\) 是 L 的元素的一个族，M 是由这些 \(x_i\) 组成的集合，U 是所有使 x 在 f 中可代入的 \(f \in \mathrm{K}((\mathrm{X}_i)_{i \in I})\) 组成的集合， \(\varphi\) 是 U 到 L 中的同态 \(f \mapsto f(\mathbf{x})\) 。则 \(\varphi(\mathrm{U})\) 是由 \(K \cup M\) 生成的 L 的子域。

设 \(L'\) 为由 \(K \cup M\) 生成的 L 的子域。我们有

\[
\mathbf {K} \cup \mathbf {M} \subset \varphi (\mathbf {U}) \subset \mathbf {L} ^ {\prime}
\]

且 \(\varphi(\mathbf{U})\) 是 \(\mathbf{L}\) 的一个子环。现在命题 2 蕴含 \(\varphi(\mathbf{U})\) 是 \(\mathbf{L}\) 的一个子域，从而 \(\varphi(\mathbf{U}) = \mathbf{L}'\) 。

设 \(f \in \mathbf{K}((\mathbf{X}_i)_{i \in I})\) ，并设 \((g_i)_{i \in I}\) 是 \(\mathbf{K}((\mathbf{Y}_l)_{l \in L})\) 的元素的一个族。若 \((g_{i})\) 在 f 中可代入，则 \(f((g_{i}))\) 是 \(\mathbf{K}((\mathbf{Y}_{l})_{l\in L})\) 的一个元素。特别地， \((\mathbf{X}_{i})_{i\in I}\) 在 f 中可代入，且 \(f=f((\mathbf{X}_{i})_{i\in I})\) 。

命题 3. --- 设 E 是 K 上的一个结合、交换、含单位元且非零的代数。设 \(f \in \mathbf{K}((\mathbf{X}_i)_{i \in \mathbb{I}})\) ，且对每个 \(i \in \mathbb{I}\) ，设 \(g_i \in \mathbf{K}((\mathbf{Y}_l)_{l \in \mathbb{I}})\) 。给定 E 的元素的一个族 \(\mathbf{y} = (y_l)_{l \in \mathbb{L}}\) ，假设 \(\mathbf{y}\) 可代入到每个 \(g_i\) 中，且 \((g_i(\mathbf{y}))_{i \in \mathbb{I}}\) 在 \(f\) 中可代入。那么：

\begin{enumerate}
\def\labelenumi{(\roman{enumi})}
\item
  \((g_{i})_{i\in I}\) 在 \(f\) 中可代入；
\item
  若以 \(h\) 表示元素 \(f((g_i))\) （ \(\mathbf{K}((\mathbf{Y}_l)_{l \in L})\) 的），则 \(\mathbf{y}\) 在 \(h\) 中可代入，且 \(h(\mathbf{y}) = f((g_i(\mathbf{y})))\) 。
\end{enumerate}

我们可以取 I 为有限集。根据假设，对每个 \(i \in I\) , \(g_i\) 可以写成 \(p_i/q_i\) 的形式，其中 \(p_i\) , \(q_i \in K[(Y_l)_{l \in L}]\) 且 \(q_i(\mathbf{y})\) 在 E 中可逆。同样， \(f\) 可以写成 \(u/v\) ，其中 \(u\) , \(v \in K[(X_i)_{i \in I}]\) 且 \(v((g_i(\mathbf{y})))\) 可逆。设 \(m = \sup(\deg u, \deg v)\) ，并设 \(w = \prod_{i \in I} q_i \in K[(Y_l)_{l \in L}]\) , \(u_1 = u((g_i))w^m\) , \(v_1 = v((g_i))w^m\) 。多项式 \(u\) 是单项式 \(\prod_{i \in I} X_i^{\nu_i}\) （满足 \(\sum_{i \in I} \nu_i \leqslant m\) ）的 K-线性组合。我们有 \(w^m \prod_{i \in I} g_i^{\nu_i} = w^m\left(\prod_{i \in I} p_i^{\nu_i}\right)\left(\prod_{i \in I} q_i^{\nu_i}\right)^{-1} \in K[(Y_l)_{l \in L}]\) （由 \(m\) 的取法）。因此 \(u_1 \in K[(Y_l)_{l \in L}]\) ，类似地 \(v_1 \in K[(Y_l)_{l \in L}]\) 。此外， \(v_1(\mathbf{y}) = (w(\mathbf{y}))^m v((g_i(\mathbf{y})))\) 可逆。因此 \(v_1 \neq 0\) （因为 \(E \neq 0\) ），从而 \(v((g_i)) \neq 0\) 。于是族 \((g_i)\) 在 \(f\) 中可代入。此外，我们有 \(f((g_i)) = u_1/v_1\) ，因此 \(\mathbf{y}\) 在 \(h = f((g_i))\) 中可代入，且 \(h(\mathbf{y}) = u_1(\mathbf{y})/v_1(\mathbf{y}) = u((g_i(\mathbf{y})))/v((g_i(\mathbf{y}))) = f((g_i(\mathbf{y})))\) 。

设 K 是一个交换域，E 是一个交换、结合且含单位元的 K-代数，并设 \(f \in \mathbf{K}((\mathbf{X}_i)_{i \in I})\) 。设 \(T_f\) 是所有 \(\mathbf{x} = (x_i)_{i \in I} \in E^{I}\) （在 \(f\) 中可代入者）组成的集合。映射 \(\mathbf{x} \mapsto f(\mathbf{x})\) （ \(T_f\) 到 E 中的）称为与 \(f\) （以及 E）相伴的有理函数；我们有时将它记作 \(\tilde{f}\) 。若 \(g \in \mathbf{K}((\mathbf{X}_i)_{i \in I})\) ，我们有 \(T_f \cap T_g \subset T_{f + g}\) , \(T_f \cap T_g \subset T_{fg}\) ，因此与 \(f + g\) （相应地 \(fg\) ）相伴的有理函数定义在 \(T_f \cap T_g\) 上，并且在该集合上与 \(\tilde{f} + \tilde{g}\) （相应地 \(\tilde{f}\tilde{g}\) ）取相同的值。设 \(T_f'\) 为所有 \(\mathbf{x} \in T_f\) （使 \(f(\mathbf{x})\) 可逆者）组成的集合；若 \(\mathbf{x} \in T_f'\) ，则 \(\mathbf{x}\) 在 \(1/f\) 中可代入，且与 \(1/f\) 相伴的有理函数在 \(\mathbf{x}\) 处取值 \(f(\mathbf{x})^{-1}\) 。

如果 K 是一个无限交换域， \(f \in \mathbf{K}((\mathbf{X}_i)_{i \in I})\) , \(g \in \mathbf{K}((\mathbf{X}_i)_{i \in I})\) ，且 \(\tilde{f}\) , \(\tilde{g}\) 是与 f, g（以及 K）相伴的有理函数，并且如果 \(\tilde{f}(\mathbf{x}) = \tilde{g}(\mathbf{x})\) （对所有 \(x \in T_f \cap T_g\) ），则 f = g。因为如果 f = u/v 且 \(g = u_1 / v_1\) ，其中 u, v, \(u_1\) , \(v_1\) 是多项式，我们有 \(u(\mathbf{x})v_1(\mathbf{x}) = u_1(\mathbf{x})v(\mathbf{x})\) （对所有满足 \(v(x)v_1(x) \neq 0\) 的 x），因此 \(uv_1 = u_1v\) （IV，第 18 页，定理 2）。因此映射 \(f \mapsto \tilde{f}\) 是单射的，我们将经常把 f 与 \(\tilde{f}\) 等同。

* 利用 \(\mathbf{K}[(\mathbf{X}_i)_{i\in I}]\) 的唯一因子分解性（交换代数，第七章，第 3 节，第 502 页及定理 2 的推论，第 506 页），容易证明以下结论：对每一个 \(f\in \mathbf{K}((\mathbf{X}_i)_{i\in I})\) ，存在 \(u,v\in \mathbf{K}[(\mathbf{X}_i)_{i\in I}]\) 使得：

\begin{enumerate}
\def\labelenumi{\arabic{enumi})}
\item
  \(f = u / v\)
\item
  \(\mathbf{x} \in \mathbf{K}^{I}\) 在 \(f\) 中可代入的充分必要条件是 \(v(\mathbf{x}) \neq 0\) 。
\end{enumerate}

